\subsection{Dimension descent and pointed-orbit restrictions for every Morita partner}
\label{subsec:focus3472-morita}

The following argument does not assume that either rank-two based module
is categorifiable.  It applies to every Morita partner of every complex
categorification of $\sha{R8_3_3472}$.

\begin{lemma}[Dimension-field descent]
\label{lem:focus-dimension-descent}
Let $\mathcal B$ be a complex fusion category such that its Drinfeld
center has no nontrivial invertible objects.  Then
\[
 \Q\bigl(\FPdim(Y):Y\in\Irr(\mathcal B)\bigr)
 =\Q\bigl(\FPdim(\mathcal B)\bigr).
\]
\end{lemma}
\begin{proof}
Write $d_Y=\FPdim(Y)$, $D=\sum_Yd_Y^2$ and $K=\Q(d_Y:Y)$.
Every conjugate of $d_Y^2$ is an eigenvalue of the positive semidefinite
integer matrix $N_YN_{Y^*}=N_YN_Y^{\mathsf T}$; consequently $K$ is
totally real.  Every embedding $\sigma:K\longrightarrow\mathbb R$
satisfies $|\sigma(d_Y)|\leq d_Y$, since $d_Y$ is the spectral radius
of the nonnegative fusion matrix $N_Y$.
If $\sigma(D)=D$, equality in
\[
 D=\sum_Y\sigma(d_Y)^2\leq\sum_Yd_Y^2=D
\]
gives $\sigma(d_Y)=\varepsilon_Yd_Y$, with
$\varepsilon_Y\in\{1,-1\}$. Applying $\sigma$ to the dimension
equation for $XY$, equality in the triangle inequality gives
$\varepsilon_Z=\varepsilon_X\varepsilon_Y$ whenever
$N_{XY}^{Z}>0$. The scalars $\varepsilon_Y$ therefore define a tensor
automorphism of the identity functor, equivalently a central object
whose underlying object is the unit.  The hypothesis forces this
automorphism to be trivial. Thus every embedding of $K$ fixing
$\Q(D)$ is the identity, and separability gives $K=\Q(D)$.
\end{proof}

\begin{proposition}[All-Morita pointed restriction]
\label{prop:focus3472-morita-invertibles}
Let $\mathcal C$ be any complex categorification of
$\sha{R8_3_3472}$ and let $\mathcal B$ be Morita equivalent to
$\mathcal C$. Put $r=\sqrt2$ and $\delta=7+5r$. Then every simple
of $\mathcal B$ has dimension $a+br$, with $a,b\in\NN$, and
\[
 \sum a^2=350,\qquad \sum b^2=175,\qquad\sum ab=245.
 \tag{\ensuremath{\dagger}}\label{eq:focus3472-moments}
\]
The integral-dimensional subcategory is pointed, and its order is
either $1$ or $7$. In the latter case the complete dimension profile
is $1^7,\delta^7$.
\end{proposition}
\begin{proof}
The source tensor has only one invertible basis element, the unit,
and $N_{ii}^{i}=2$ for each of its seven nonunit basis elements.
An invertible object of $Z(\mathcal C)$ must therefore forget to the
unit. Its half-braiding is a tensor automorphism of the identity,
with nonzero scalars $\lambda_i$ satisfying
$\lambda_i^2=\lambda_i$ because $X_i$ occurs in $X_i^2$.
All $\lambda_i$ are one. Thus $Z(\mathcal C)$ has no nontrivial
invertibles, and neither does $Z(\mathcal B)$ by Morita invariance.

Morita invariance also gives
$D=\FPdim(\mathcal B)=700+490r$. The preceding lemma puts every
$d_Y$ in $\Q(r)$, whose ring of integers is $\Z[r]$.
Write $d_Y=a_Y+b_Yr$. Positivity of $d_Y$ and the bound
$|a_Y-b_Yr|\leq a_Y+b_Yr$ imply $a_Y,b_Y\geq0$.
Orthogonality of the two distinct dimension characters gives
$\sum(a_Y^2-2b_Y^2)=0$. Combining this with $\sum d_Y^2=D$
proves \eqref{eq:focus3472-moments}.

The integral-dimensional simples form a fusion subcategory because
all radical coefficients are nonnegative. Let $h$ be its dimension.
The integral-dimension Galois Gram bound gives $h\leq7$, since
$\operatorname{Tr}_{\Q(r)/\Q}(D^{-1})=1/7$.
Categorical Lagrange integrality of $D/h$ gives
$1400/h,9800/h^2\in\Z$, whence $h\in\{1,2,5,7\}$.
This subcategory is pointed. The only nonpointed integral profiles
at these dimensions are $1,2$ at dimension five and $1,1,1,2$ at
dimension seven. Their unique dimension-two simple is self-dual and
fixed by every invertible, giving respectively the impossible
dimension equations $4=1+2m$ and $4=3+2m$.
Write $g=h$ for the number of invertibles.

We use strong Lagrange divisibility: if $\mathcal D\subseteq\mathcal B$
is a fusion subcategory and $\mathcal M$ is an indecomposable
$\mathcal D$-module component of $\mathcal B$, then
$\FPdim(\mathcal M)/\FPdim(\mathcal D)$ is an algebraic integer
\cite[Proposition~5.7]{Integral}. For the pointed subcategory and a
left orbit $\mathcal O$ of simples of common dimension $d$, this says
$|\mathcal O|d^2/g$ is integral.
If $g=2$, a fixed simple of dimension $a+br$ must have
$d^2/2$ integral, so $a$ is even. Free orbits contribute an even
integer to $\sum ab$, as do fixed simples. This contradicts
$\sum ab=245$.

Suppose $g=5$. A fixed simple has $d^2/5$ integral, so
$a^2+2b^2\equiv 2ab\equiv0\pmod5$, implying $5\mid a,b$.
From \eqref{eq:focus3472-moments}, after removing the invertibles,
\[
 \sum(5a-7b)^2=50.
\]
Together with $a^2\leq345$, $b^2\leq175$ this leaves only
$(a,b)=(15,10)$ for a fixed nonintegral simple. There can be at most
one. After removing it, the remaining nonintegral dimensions would
satisfy $\sum a^2=120$, $\sum b^2=75$, $\sum ab=95$, contrary to
$95^2>120\cdot75$. Hence every nonintegral orbit has length five.
For orbit representatives, the moment equations become
\[
 \sum a^2=69,\qquad\sum b^2=35,\qquad\sum ab=49,
 \qquad\sum(5a-7b)^2=10.
\]
The complete possible types are
$(1,1),(2,1),(3,2),(4,3),(5,4),(6,4),(7,5)$.
Their nonnegative multiplicities give exactly the following four
lists of nonintegral orbit-representative dimensions:
\[
\begin{array}{c|l}
 \mathrm A&2+r,\ 4+3r,\ 7+5r\\
 \mathrm B&1+r,\ (4+3r)^2,\ 6+4r\\
 \mathrm C&1+r,\ (3+2r)^4,\ (4+3r)^2\\
 \mathrm D&(1+r)^2,\ (3+2r)^2,\ 7+5r.
\end{array}
\]
Superscripts denote multiplicities of free orbits. This finite
seven-type enumeration is replayed with exact integer arithmetic.
In A, the product $XX^*$ for $\FPdim(X)=2+r$ contains exactly one
invertible, since its stabilizer is trivial. Its nonintegral part
would have dimension $5+4r$, which is no nonnegative combination of
the three available dimensions (their rational parts are even or
at least seven). Thus A is impossible.

For B and C, put $u=1+r$. There is one free orbit of $u$-dimensional
simples, necessarily stable under duality. The right orbit of $X^*$
is free, since its stabilizer is dual to the left stabilizer of $X$,
and therefore fills this same five-element orbit. Hence for any two
$u$-dimensional simples $X,Y$ there is exactly one $g\in C_5$ with
$Y\simeq X^*g$. Their product contains exactly one invertible
and has nonintegral part $u^2-1=2u$. No larger available dimension
fits that remainder. They therefore generate, with the five
invertibles, a fusion subcategory of profile $1^5,u^5$.

In D, there are two $u$-orbits, two $u^2$-orbits and one $u^3$-orbit.
Right multiplication by $C_5$ preserves each left orbit: it acts on
the left orbits of any fixed dimension, and there are fewer than
five of these. Hence identifying the five invertibles with the
unit gives a based fusion-ring quotient on the six left orbits,
of dimensions $1,u,u,u^2,u^2,u^3$. Here are the details that ensure
this is a based quotient, rather than merely a dimension quotient.
For orbit representatives $x,y,z$, set
$\overline N_{[x][y]}^{[z]}=\sum_{g\in C_5}N_{xy}^{gz}$.
Left translation of $x$ translates the product, and replacing $y$
by $gy$ does the same because $xg$ belongs to the left orbit of $x$.
The coefficients are therefore independent of representatives and
inherit associativity and duality. Their unit coefficient counts
$g$ with $y\simeq x^*g$, so it is one precisely for dual orbits.
Finally, Frobenius reciprocity follows from
\[
 \sum_gN_{xy}^{gz}
 =\sum_gN_{y,z^*g^{-1}}^{x^*}
 =\sum_gN_{y,z^*}^{x^*g}
 =\overline N_{[y][z^*]}^{[x^*]},
\]
since right multiplication runs through the same free left orbit.
Write $A,B$ for the two $u$-dimensional basis elements of this quotient.
A product with a unit term has the form $1+$ two $u$-terms;
without a unit term it is a single $u^2$-dimensional basis element.
If $A^*=B$, the squares $A^2,B^2$ have no $u$-terms. Cyclic
Frobenius reciprocity then gives
$[AB:A]=[BB:B]=0$ and $[AB:B]=[AA:A]=0$, contradicting
$\FPdim(AB-1)=2u$. Thus both are self-dual. Now $AB$ is one
$u^2$-dimensional basis element, so reciprocity gives
$[A^2:B]=[AB:A]=0$ and similarly $[B^2:A]=0$. Consequently
$A^2=1+2A$ and $B^2=1+2B$. Lifting either orbit back to
$\mathcal B$ again gives a fusion subcategory of profile $1^5,u^5$.
Lemma~\ref{lem:focus-c5-silver} below excludes this profile and
finishes the order-five case.

Finally, if $g=7$, equality in the Galois Gram bound forces
$5a=7b$ for every nonintegral dimension. Thus those dimensions are
$k\delta$, with positive integers $k$ satisfying $\sum k^2=7$.
The possibilities are $\delta^7$ and $\delta^3,(2\delta)^1$.
Since $\delta^2=1+14\delta$, a $\delta$-dimensional simple has
trivial stabilizer among the seven invertibles. Its orbit has seven
members, excluding the second possibility.
\end{proof}

\begin{lemma}[The order-five silver-ratio profile]
\label{lem:focus-c5-silver}
There is no complex fusion category with dimension profile
$1^5,(1+\sqrt2)^5$ whose invertibles act freely on the
noninvertible simples.
\end{lemma}
\begin{proof}
Let $g$ generate the invertibles and choose a noninvertible $X$.
The left action is free and transitive. The right action is too:
the coefficients of invertibles in $XX^*$ and $X^*X$ are each
forced by $\FPdim(X)^2=1+2\FPdim(X)$ to sum to one. Thus
$Xg=g^aX$ for some $a\in(\Z/5\Z)^\times$.
Since the dual remains in the same free orbit, dualizing the
relation gives $a^2=1$, so $a=\pm1$.
One can choose $X=X^*$: for $a=1$ shift its index by half its
duality translation, and for $a=-1$ the involution forces that
translation to vanish. Write
\[
 X^2=1+\sum_{j\in\Z/5\Z}m_jg^jX,
 \qquad m_j\in\NN,\qquad\sum m_j=2.
\]
If $a=-1$, comparing $X^2g$ with $X(Xg)$ forces $m_j=m_{j+2}$,
which would make all five $m_j$ equal, an impossibility.
If $a=1$, duality gives $m_j=m_{-j}$. There are two possibilities
up to changing the generator: $m_0=2$, or $m_1=m_{-1}=1$.
The first contains the rank-two fusion subcategory
$X^2=1+2X$, excluded by the rank-two classification
\cite{Ostrik2}.

In the second case the ring is commutative. Choose a primitive fifth
root $\zeta$ and a one-dimensional character with
$\chi(g)=\zeta$, $\chi(X)=x$, where
$x^2-(\zeta+\zeta^{-1})x-1=0$. Its formal codegree is
$f=5(1+x^2)$. Since $t=\zeta+\zeta^{-1}$ satisfies
$t^2+t-1=0$, eliminating $x,t$ gives
\[
 f^4-55f^3+1000f^2-7250f+18125=0.
\]
The polynomial is irreducible: modulo two it is $z^4+z^3+1$,
which has neither a root nor the factor $z^2+z+1$.
The formal-codegree $d$-number criterion
\cite[Theorem~1.2 and Lemma~2.7]{OstrikCode} would
require $18125\mid55^4$, whereas the remainder is $15625$.
This excludes the second case without a pivotality assumption.
\end{proof}

\begin{corollary}
The first rank-two branch of
Proposition~\ref{prop:creative-r3472-three} retains only four of its
seven dimension spectra: its three one-invertible spectra and the
seven-invertible spectrum $1^7,(7+5\sqrt2)^7$.
\end{corollary}
\begin{proof}
The all-Morita proposition removes the remaining order-two spectrum
and both order-five spectra. This is still a conditional restriction
on the rank-two construction route: neither its lift nor a
categorification of the source ring is asserted.
\end{proof}

The verifier in
\path{evidence/r3472_focus_20261022/morita/verify.py} checks the
source tensor, all arithmetic bounds, the complete four-profile
enumeration and the irreducible formal-codegree obstruction.
The general categorification status of $\sha{R8_3_3472}$ remains open.
